\documentclass{scrartcl}
\usepackage{graphicx}  % required for inserting images

% for cyrillic and Bulgarian language
\usepackage[utf8]{inputenc}
\usepackage[T2A]{fontenc}
\usepackage[bulgarian]{babel}

% math packages
\usepackage{amsmath}
\usepackage{float}
\usepackage{amsfonts}
\usepackage{amssymb}
\usepackage{amsthm}
\usepackage{cancel}

\usepackage{ragged2e}  % adds FlushLeft

\usepackage{listings}
\usepackage{xcolor}

\definecolor{codegreen}{rgb}{0,0.6,0}
\definecolor{codegray}{rgb}{0.5,0.5,0.5}
\definecolor{codepurple}{rgb}{0.58,0,0.82}
\definecolor{backcolour}{rgb}{0.95,0.95,0.92}

\lstdefinestyle{mystyle}{
    language=Octave,
    backgroundcolor=\color{backcolour},   
    commentstyle=\color{codegreen},
    keywordstyle=\color{magenta},
    numberstyle=\tiny\color{codegray},
    stringstyle=\color{codepurple},
    basicstyle=\ttfamily\footnotesize,
    breakatwhitespace=false,         
    breaklines=true,
    captionpos=b,
    keepspaces=false,                 
    numbersep=5pt,
    showspaces=false,
    showstringspaces=false,
    showtabs=false,
    tabsize=2,
    columns=fullflexible,
}

\lstset{style=mystyle}

\usepackage[colorlinks=true, linkcolor=blue, urlcolor=cyan]{hyperref}

\newtheorem{theorem}{Теорема}
\newtheorem{definition}{Дефиниция}
\newtheorem{statement}{Твърдение}

\title{Диференциални уравнения - линейни уравнения}
\author{Ивайло Андреев}
\date{15 май 2026}

\begin{document}
\maketitle

\tableofcontents

\newpage

\begin{FlushLeft}

В цялата тема важи, че $\Delta$ е област, тоест отворено и свързано множество.

\section{Линейно диференциалното уравнение}

\begin{definition}
(Комплекснозначна функция)

Комплекснозначна функция наричаме функция от вида:

$$g(t) = u(t) + \mathrm{i}v(t)$$

където $u: \Delta \rightarrow \mathbb{R}$ е реалната част на $g$ и $v: \Delta \rightarrow \mathbb{R}$ е имагинерната част на $g$.
\end{definition}

\begin{definition}
(Непрекъсната комплекснозначна функция)

Нека $g(t) = u(t) + \mathrm{i}v(t)$ е комплекснозначна функция.

Тогава $g(t)$ е непрекъсната, ако функциите $u(t)$ и $v(t)$ са непрекъснати в $\Delta$.
\end{definition}

\begin{definition}
(Линейно уравнение от ред $n$)

Нека $y(t)$ е $n$ пъти диференцируема, тоест $y \in C^n(\Delta)$.

Линейно обикновено диференциално уравнение от ред $n$ за $y(t)$ наричаме уравнение от вида:

\begin{equation}\label{eq:general_linear}
    \sum_{i=0}^n a_i(t)y^{(i)} = f(t)
\end{equation}

където $a_0(t), a_1(t), \dots, a_n(t)$ и $f(t)$ са непрекъснати комплекснозначни функции и $a_n(t) \ne 0 \; \forall t \in \Delta$.
\end{definition}

Можем без ограничение на общността да считаме, че $a_n(t) \equiv 1$, защото $a_n(t) \ne 0 \; \forall t \in \Delta$ и можем да разделим на $a_n(t)$.

\section{Теорема за съществуване и единственост}

\begin{theorem}\label{th:unique_sol}
(Теорема за съществуване и единственост)

Нека е дадено линейното уравнение \eqref{eq:general_linear}.
Нека $t_0 \in \Delta$ и $\eta_0, \eta_1, \dots, \eta_{n-1} \in \mathbb{C}$.

Тогава следната задача на Коши:

$$
\begin{cases}
\displaystyle \sum_{i=0}^n a_i(t)y^{(i)} = f(t) \\
y(t_0) = \eta_0 \\
y'(t_0) = \eta_1 \\
\vdots \\
y^{(n-1)}(t_0) = \eta_{n-1} \\
\end{cases}
$$

има единствено решение и то е дефинирано в целия интервал $\Delta$.
\end{theorem}

\begin{definition}
(Хомогенно линейно уравнение от ред $n$)

Наричаме уравнението \eqref{eq:general_linear} хомогенно, ако $f(t) \equiv 0$. Иначе го наричаме нехомогенно.
\end{definition}

Хомогенното линейно уравнение има вида

\begin{equation}\label{eq:homog}
    \sum_{i=0}^n a_i(t)y^{(i)} = 0
\end{equation}

\begin{definition}
(Уравнение с постоянни коефициенти)

Казваме, че уравнението \eqref{eq:general_linear} е с постоянни коефициенти, ако функциите $a_0(t), a_1(t), \dots, a_n(t)$ са константни.
\end{definition}

\section{Критерий за линейна независимост на решения}

\begin{definition}
(Линейна независимост на система от функции)

Системата от функции $y_1(t), \dots, y_n(t)$ се нарича линейно независима в $\Delta$, ако от това, че

$$\sum_{j=1}^n C_jy_j(t) = 0 \quad \forall t \in \Delta$$

следва, че

$$(C_1, \dots, C_n) = (0, \dots, 0)$$
\end{definition}

\begin{definition}
(Линейна зависимост на система от функции)

Системата от функции $y_1(t), \dots, y_n(t)$ се нарича линейно зависима в $\Delta$, ако

$$\sum_{j=1}^n C_jy_j(t) = 0 \quad \forall t \in \Delta$$

и

$$(C_1, \dots, C_n) \ne (0, \dots, 0)$$
\end{definition}

\begin{definition}
(Детерминанта на Вронски)

Нека $y_1, \dots, y_n \in C^{(n-1)}(\Delta)$.

Детерминанта на Вронски за функциите $y_1, \dots, y_n$ наричаме:

$$
W(y_1, \dots, y_n)(t) = 
\begin{vmatrix}
y_1(t) & y_2(t) & \dots & y_n(t) \\
y_1'(t) & y_2'(t) & \dots & y_n'(t) \\
\vdots & \vdots & \ddots & \vdots \\
y_1^{(n-1)}(t) & y_2^{(n-1)}(t) & \dots & y_n^{(n-1)}(t)
\end{vmatrix}
$$
\end{definition}

Стойността на детерминантата на Вронски в точка $t_0$ ще бележим с $W(t_0)$.

\begin{theorem}
(Критерий за линейна независимост на система от функции)

Нека $y_1, \dots, y_n \in C^{(n-1)}(\Delta)$ са решение на хомогенното линейно уравнение \eqref{eq:homog} и нека $W$ е тяхната детерминанта на Вронски.

Тогава ако $\exists t_0 \in \Delta : W(t_0) \ne 0$, то системата от функции $y_1, \dots, y_n$ е линейно независима в $\Delta$.
\end{theorem}

\textbf{Доказателство:}

Нека $C_1, \dots, C_n \in \mathbb{C}$ са такива константи, за които е изпълнено следното равенство:

$$\sum_{j=1}^n C_jy_j(t) = 0$$

Ще покажем, че $(C_1, \dots, C_n) = (0, \dots, 0)$, от което ще следва, че системата от функциите е линейно независима в $\Delta$.

Диференцираме последователно $n-1$ пъти и получаваме следната система.

$$
\begin{cases}
    C_1y_1(t) + \dots + C_ny_n(t) = 0\\
    C_1y_1'(t) + \dots + C_ny_n'(t) = 0\\
    \vdots\\
    C_1y_1^{(n-1)}(t) + \dots + C_ny_n^{(n-1)}(t) = 0\\
\end{cases}
$$

При $t = t_0$ получаваме следното матрично уравнение:

$$
\begin{pmatrix}
    y_1(t_0) & y_2(t_0) & \dots & y_n(t_0) \\
    y_1'(t_0) & y_2'(t_0) & \dots & y_n'(t_0) \\
    \vdots & \vdots & \ddots & \vdots \\
    y_1^{(n-1)}(t_0) & y_2^{(n-1)}(t_0) & \dots & y_n^{(n-1)}(t_0)
\end{pmatrix}
\begin{pmatrix}
    C_1\\
    C_2\\
    \vdots\\
    C_n
\end{pmatrix}
=
\begin{pmatrix}
    0\\
    0\\
    \vdots\\
    0
\end{pmatrix}
$$

Детерминантата на матрицата вляво е детерминантата на Вронски, която по условие е ненулева в точката $t_0$ ($W(t_0) \ne 0$).

Така системата има единствено решение при $(C_1, \dots, C_n) = (0, \dots, 0)$, от което следва, че дадената система от функции е линейно независима в $\Delta$.

\begin{theorem}
(Критерий за линейна зависимост на система от функции)

Нека $y_1, \dots, y_n \in C^{(n-1)}(\Delta)$ са решение на хомогенното линейно уравнение \eqref{eq:homog} и нека $W$ е тяхната детерминанта на Вронски.

Тогава ако $\exists t_0 \in \Delta : W(t_0) = 0$, то системата от функции $y_1, \dots, y_n$ е линейно зависима в $\Delta$.
\end{theorem}

\section{Линейно пространство}

\begin{definition}
(Фундаментална система решения)

Решенията на хомогенното линейно уравнение \eqref{eq:homog} наричаме фундаментална система решения (ФСР) в $\Delta$, ако те са линейно независими.
\end{definition}

\begin{theorem}
(Критерий за фундаментална система решения)

Нека $y_1(t), \dots, y_n(t)$ са решения на хомогенното уравнение \eqref{eq:homog} в $\Delta$ и нека $W$ е тяхната детерминанта на Вронски.

Тогава следните твърдения са еквивалентни:

\begin{enumerate}
    \item $\exists t_0 \in \Delta : W(t_0) \ne 0$
    \item $y_1(t), \dots, y_n(t)$ образуват ФСР в $\Delta$ за уравнението
    \item $W(t) \ne 0 \; \forall t \in \Delta$
\end{enumerate}
\end{theorem}

\begin{theorem}
(Линейно пространство от решения)

Множеството от всички решения на хомогенното линейно уравнение \eqref{eq:homog} образува линейно пространство с размерност $n$.
\end{theorem}

\textbf{Доказателство:}

Нека означим множеството от всички решения на хомогенното уравнение с $M$.

Ще покажем, че $M$ е изоморфно с линейно пространство $\mathbb{C}^n$.

За да говорим за изоморфизъм е необходимо множеството $M$ да е линейно пространство.

\underline{0. $M$ е линейно пространство}

Нека $y_1(t)$ и $y_2(t)$ са решения на хомогенното уравнение и нека $\alpha, \beta \in \mathbb{C}$.

Ще покажем, че $\alpha y_1(t) + \beta y_2(t)$ също е решение.

\begin{align*}
    \sum_{i=0}^n a_i(\alpha y_1 + \beta y_2)^{(i)} & = \sum_{i=0}^n a_i(\alpha y_1^{(i)} + \beta y_2^{(i)})\\
    & = \sum_{i=0}^n a_i\alpha y_1^{(i)} + \sum_{i=0}^n a_i \beta y_2^{(i)}\\
    & = \alpha \sum_{i=0}^n a_i y_1^{(i)} + \beta \sum_{i=0}^n a_i y_2^{(i)}\\
    & = \alpha \times 0 + \beta \times 0 = 0
\end{align*}

Показахме, че $M$ е линейно пространство. Остава да покажем, че размерността му е равна на $n$. Ще го покажем, като построим биективно линейно изображение (изоморфизъм).

Нека $H$ е множеството от всички наредени $n$-торки с комплексни числа, които представляват начални условия за линейното уравнение от $n$-ти ред. Нека въведем изображението

$$\psi: M \to H$$

То съпоставя на всяко решение началните му стойности в точката $t_0 \in \Delta$.

Нека $y_0(t)$ е решение на задачата на Коши за уравнението с начални условия $(\eta_0, \dots, \eta_{n-1})$ в точката $t_0$. Тогава

$$\psi(y_0(t)) = (y(t_0), \dots, y^{(n-1)}(t_0)) = (\eta_0, \dots, \eta_{n-1})$$

\underline{1. $\psi$ е линейно изображение}

Нека $y_3(t)$ и $y_4(t)$ са решения на хомогенното уравнение и нека $\gamma, \rho \in \mathbb{C}$.

\begin{align*}
    \psi(\gamma y_3 + \rho y_4) & = (\gamma y_3(t_0) + \rho y_4(t_0), \dots, \gamma y_3^{(n-1)}(t_0) + \rho y_4^{(n-1)}(t_0))\\
    & = (\gamma y_3(t_0), \dots, \gamma y_3^{(n-1)}(t_0)) + (\rho y_4(t_0), \dots, \rho y_4^{(n-1)}(t_0))\\
    & = \gamma(y_3(t_0), \dots, y_3^{(n-1)}(t_0)) + \rho (y_4(t_0), \dots, y_4^{(n-1)}(t_0))\\
    & = \gamma \psi(y_3) + \rho \psi(y_4)
\end{align*}

\underline{2. $\psi$ е сюрекция}

От теоремата за съществуване и единственост \ref{th:unique_sol} следва, че за всяка $n$-торка от начални условия съществува единствено решение. Следователно $\psi$ е сюрекция.

\underline{3. $\psi$ е инекция}

От теоремата за съществуване и единственост \ref{th:unique_sol} следва, че за $n$-торката от начални условия $(0, \dots, 0)$ съществува единствено решение, което е нулевото. Оттук $\text{Ker }\psi = \{0\}$. Следователно $\psi$ е инекция.

Тъй като $H \cong \mathbb{C}^n$, то $\dim H = n$.
Понеже $M \cong H$, получаваме $\dim M = n$.

\begin{statement}\label{st:general_solution}
(Общо решение на хомогенно уравнение)

Нека е дадено хомогенно уравнение от ред $n$ от вида \eqref{eq:homog}.
Нека фундаменталната система решения за хомогенното уравнение има вида:
$$\text{ФСР} = \{\varphi_1(t), \dots, \varphi_n(t)\}$$

Тогава хомогенното уравнение има общо решение от вида:

$$y(t) = y_0(t) = \sum_{j=1}^n C_j \varphi_j(t)$$

където $C_1, \dots, C_n \in \mathbb{C}$.
\end{statement}

\section{Структура на решения на нехомогенно уравнение}

\begin{definition}
(Квазиполином)

Нека $\alpha, \beta \in \mathbb{R}$. Нека $k \in \mathbb{N}$ и нека $P_k(t)$ е полином от $k$-та степен.

Тогава квазиполином наричаме функция от един от следните видове:

\begin{enumerate}
    \item $P_k(t)e^{\alpha t}$
    \item $P_k(t)e^{\alpha t}\cos{(\beta t)}$
    \item $P_k(t)e^{\alpha t}\sin{(\beta t)}$
\end{enumerate}
\end{definition}

\begin{statement}
(Частно решение - квазиполином)

Нека е дадено нехомогенно уравнение от вида \eqref{eq:general_linear} с \underline{постоянни коефициенти} и нека $f(t)$ е квазиполином.

Тогава нехомогенното уравнение има частно решение във вид на квазиполином.
\end{statement}

\begin{statement}
(Частно решение - метод на Лагранж)

Нека е дадено нехомогенно уравнение от вида \eqref{eq:general_linear} и нека $f(t)$ е произволна непрекъсната функция.
Нека фундаменталната система решения за \underline{хомогенното} уравнение при $f(t) \equiv 0$ има вида:
$$\text{ФСР} = \{\varphi_1(t), \dots, \varphi_n(t)\}$$

Тогава нехомогенното уравнение има частно решение от вида:

$$y_1(t) = \sum_{j=1}^n b_j(t) \varphi_j(t)$$

където $b_1(t), \dots, b_n(t) : \Delta \rightarrow \mathbb{C}$ са непрекъснати функции, които удовлетворяват условията:

$$
\begin{cases}
    \displaystyle\sum_{j=1}^n b_j'(t)\varphi_j(t) = 0\\
    \displaystyle\sum_{j=1}^n b_j'(t)\varphi_j'(t) = 0\\
    \vdots\\
    \displaystyle\sum_{j=1}^n b_j'(t)\varphi_j^{(n-2)}(t) = 0\\
    \displaystyle\sum_{j=1}^n b_j'(t)\varphi_j^{(n-1)}(t) = \dfrac{f(t)}{a_n(t)}\\
\end{cases}
$$
\end{statement}

\begin{theorem}
(Структура на решението на нехомогенно уравнение)

Нека е дадено нехомогенно уравнение от вида \eqref{eq:general_linear}.
Нека $y_0(t)$ е общо решение за \underline{хомогенното} уравнение при $f(t) \equiv 0$ и нека $y_1(t)$ е частно решение на нехомогенното уравнение.

Тогава общото решение на нехомогенното уравнение има вида:

$$y(t) = y_0(t) + y_1(t)$$    
\end{theorem}

\section{Алгоритъм за намиране на общо решение на хомогенно уравнение с постоянни коефициенти}

\subsection{Алгоритъм}

Нека е дадено уравнението \eqref{eq:homog} с постоянни реални коефициенти и нека $y(t) \in \mathbb{R}$.

$$\sum_{i=0}^n a_i y^{(i)} = 0$$

На даденото линейно уравнение съответства следният характеристичен полином:

$$P(\lambda) = \sum_{i=0}^n a_i \lambda^i$$

Приравняваме характеристичния полином на $0$.

$$\sum_{i=0}^n a_i \lambda^{i} = 0$$

Получаваме полиномиално уравнение. Тъй като $a_n \ne 0$, характеристичният полином е от степен $n$. Според основната теорема на алгебрата уравнението притежава $n$ на брой комплексни корена, считани с тяхната кратност.

На всеки корен на характеристичното уравнение съответстват решения по следните правила:

\begin{enumerate}
    \item $\lambda_i \in \mathbb{R}$ е еднократен реален корен:
    
    $$\varphi_{\lambda_i}(t) = e^{\lambda_i t}$$
    
    \item $\lambda_{i} \in \mathbb{R}$ е $k$-кратен реален корен:

    $$
    \begin{cases}
        \varphi_{\lambda_{i_1}}(t) = e^{\lambda_i t}\\
        \varphi_{\lambda_{i_2}}(t) = te^{\lambda_i t}\\
        \vdots\\
        \varphi_{\lambda_{i_k}}(t) = t^{k-1}e^{\lambda_i t}
    \end{cases}
    $$
    
    \item $\lambda_{i_1, i_2} = \alpha \pm \mathrm{i} \beta\in \mathbb{C}$ е еднократен комплексно спрегнат корен:

    $$
    \begin{cases}
        \varphi_{\lambda_{i_1}}(t) = e^{\alpha t}\cos{(\beta t)}\\
        \varphi_{\lambda_{i_2}}(t) = e^{\alpha t}\sin{(\beta t)}
    \end{cases}
    $$
    
    \item $\lambda_{i_1, i_2} = \alpha \pm \mathrm{i} \beta\in \mathbb{C}$ е $k$-кратен комплексно спрегнат корен:

    $$
    \begin{cases}
        \begin{cases}
            \varphi_{\lambda_{i_{1, 1}}}(t) = e^{\alpha t}\cos{(\beta t)}\\
            \varphi_{\lambda_{i_{1, 2}}}(t) = e^{\alpha t}\sin{(\beta t)}
        \end{cases}\\
        \begin{cases}
            \varphi_{\lambda_{i_{2, 1}}}(t) = te^{\alpha t}\cos{(\beta t)}\\
            \varphi_{\lambda_{i_{2, 2}}}(t) = te^{\alpha t}\sin{(\beta t)}
        \end{cases}\\
        \vdots\\
        \begin{cases}
            \varphi_{\lambda_{i_{k, 1}}}(t) = t^{k-1}e^{\alpha t}\cos{(\beta t)}\\
            \varphi_{\lambda_{i_{k, 2}}}(t) = t^{k-1}e^{\alpha t}\sin{(\beta t)}
        \end{cases}
    \end{cases}
    $$
\end{enumerate}

След прилагане на тези правила формираме фундаментална система решения от базисните функции:

$$\text{ФСР} = \{\varphi_1(t), \dots, \varphi_n(t)\}$$

Според твърдението \ref{st:general_solution} общото решение има вида:

$$y(t) = \sum_{j=1}^n C_j \varphi_j(t)$$

където $C_1, \dots, C_n \in \mathbb{R}$.

\subsection{Пример}

$$y'' - 4y' + 4y = 0$$

На даденото линейно уравнение съответства следният характеристичен полином:

$$P(\lambda) = \lambda^2 -4 \lambda + 4$$

Приравняваме характеристичния полином на $0$.

$$\lambda^2 - 4\lambda + 4 = 0$$

$$(\lambda - 2)^2 = 0$$

Получаваме един двоен реален корен $\lambda_{1,2} = 2$.

$$\text{ФСР }=\lbrace \mathrm{e}^{2t}, \, t\mathrm{e}^{2t} \rbrace$$

Общото решение на хомогенното уравнение $y_0$ е:

$$y = C_1 \mathrm{e}^{2t} + C_2 t\mathrm{e}^{2t}, \quad C_1, \space C_2 \in \mathbb{R}$$

\end{FlushLeft}

\end{document}